\chapter{Instrumento de Coleta de Dados: Questionário de Experiência do Jogador e Usabilidade}
\label{apendice:questionario_jogadores}

Este apêndice reproduz integralmente o instrumento de coleta de dados psicométricos, investigativos e de usabilidade concebido e aplicado no protocolo de validação ecológica com jogadores humanos (Seção \ref{sec:met_validacao_humana}). O questionário foi estruturado em blocos modulares para aplicação imediata após a conclusão do inquérito interativo com o suspeito Alberto Castelo no ambiente Unity.

\section*{Instruções aos Participantes}
\textit{Caro(a) participante, agradecemos sua colaboração voluntária nesta pesquisa desenvolvida na Universidade do Estado do Amazonas (EST/UEA). O objetivo deste questionário é avaliar a sua experiência ao interrogar um suspeito controlado por Inteligência Artificial em um jogo de investigação policial. Não existem respostas certas ou erradas; o que nos interessa é exclusivamente a sua percepção sincera sobre a interação, a história, o ritmo do diálogo e o comportamento do personagem. Suas respostas são anônimas, confidenciais e voltadas estritamente a fins acadêmicos.}

\vspace{0.05cm}

% ====================================================================
\section{Bloco 0: Perfil do Participante e Hábitos Tecnológicos}
\label{sec:apendice_d_bloco0}
% ====================================================================

\noindent\textbf{Justificativa e Embasamento Científico:} A literatura de Interação Humano-Computador (IHC) e agentes conversacionais estabelece que a experiência prévia com inteligência artificial generativa e a familiaridade com jogos narrativos atuam como variáveis moderadoras cruciais na percepção de qualidade do sistema \cite{gnewuch2018faster,nass2000machines}. Jogadores com alta vivência em RPGs e jogos de dedução possuem repertório cognitivo prévio para avaliar a consistência de pistas policiais \cite{fernandez2014introduction}, enquanto o domínio prévio de LLMs calibra as expectativas em relação a respostas sintéticas ou tempos de processamento. A caracterização deste bloco assegura a correta estratificação amostral e a validade interna do estudo.

\begin{enumerate}[label=\textbf{P\arabic*.}, noitemsep, topsep=1pt, parsep=0pt]
    \item \textbf{Código de Identificação do Participante:} [\ \underline{\hspace{4cm}}\ ]
    \item \textbf{Faixa Etária:}
    \begin{itemize}[label={$\square$}, noitemsep, topsep=1pt, parsep=0pt]
        \item Menos de 18 anos
        \item 18 a 22 anos
        \item 23 a 27 anos
        \item 28 anos ou mais
    \end{itemize}
    \item \textbf{Frequência habitual com que joga jogos digitais (qualquer plataforma):}
    \begin{itemize}[label={$\square$}, noitemsep, topsep=1pt, parsep=0pt]
        \item Raramente ou nunca (menos de 1 hora semanal)
        \item Ocasionalmente (1 a 5 horas semanais)
        \item Moderadamente (5 a 14 horas semanais)
        \item Frequente / Entusiasta (mais de 14 horas semanais)
    \end{itemize}
    \item \textbf{Familiaridade com jogos focados em narrativa, investigação, RPG ou aventura (ex.: \textit{L.A. Noire, Disco Elysium, Phoenix Wright, Ace Attorney, Danganronpa, RPGs de mesa}):}
    \begin{itemize}[label={$\square$}, noitemsep, topsep=1pt, parsep=0pt]
        \item 1 --- Nenhuma experiência
        \item 2 --- Pouca experiência
        \item 3 --- Experiência moderada
        \item 4 --- Boa experiência
        \item 5 --- Ampla experiência / Jogador assíduo do gênero
    \end{itemize}
    \item \textbf{Grau de familiaridade e uso de assistentes de IA Generativa (ex.: \textit{ChatGPT, Claude, Gemini, Copilot}):}
    \begin{itemize}[label={$\square$}, noitemsep, topsep=1pt, parsep=0pt]
        \item Nunca utilizei
        \item Uso esporádico (apenas consultas ocasionais)
        \item Uso intermediário (utilizo com frequência semanal para estudos ou lazer)
        \item Uso avançado / Técnico (compreendo conceitos de prompts, contexto ou arquiteturas de LLM)
    \end{itemize}
\end{enumerate}

\vspace{0.05cm}

% ====================================================================
\section{Bloco 1: Dedução Investigativa e Veredito Narrativo}
\label{sec:apendice_d_bloco1}
% ====================================================================

\noindent\textbf{Justificativa e Embasamento Científico:} Conforme fundamentado por Ryan \cite{ryan2006interactive} e Riedl e Bulitko \cite{riedl2013interactive}, uma narrativa interativa de mistério atinge excelência quando o jogador consegue correlacionar evidências dispersas e chegar a conclusões dedutivas fundamentadas logicamente. As perguntas deste bloco visam mensurar se a arquitetura de \textit{Agentic RAG} foi capaz de transmitir os fatos canônicos do caso (álibi no jardim, dívidas com agiotas às 09:00 e o flagrante de Silas às 11:30) com clareza diegética suficiente para permitir que o participante decida a inocência ou culpa do suspeito com elevado grau de certeza, atestando a integridade informativa da experiência.

\begin{enumerate}[label=\textbf{D\arabic*.}, noitemsep, topsep=1pt, parsep=0pt]
    \item \textbf{Com base nas informações e no comportamento demonstrado por Alberto Castelo durante o interrogatório, qual é o seu veredito?}
    \begin{itemize}[label={$\square$}, noitemsep, topsep=1pt, parsep=0pt]
        \item \textbf{INOCENTE} do homicídio de Arthur Castelo
        \item \textbf{CULPADO} pelo homicídio de Arthur Castelo
    \end{itemize}
    \item \textbf{Qual o seu grau de certeza no veredito formalizado?}
    \begin{itemize}[label={$\square$}, noitemsep, topsep=1pt, parsep=0pt]
        \item 1 --- Chute / Muito incerto
        \item 2 --- Pouco convicto
        \item 3 --- Razoavelmente seguro
        \item 4 --- Convicto com base em indícios
        \item 5 --- Certeza absoluta sustentada pelas contradições do depoimento
    \end{itemize}
    \item \textbf{Quais declarações, contradições de álibi ou pistas foram determinantes para a sua conclusão?}
    
    
    
    \item \textbf{Você teve a sensação de que o suspeito tentou deliberadamente ocultar algum fato relevante ou desviar o foco das suas perguntas?}
    
    
\end{enumerate}

\vspace{0.05cm}

% ====================================================================
\section{Bloco 2: Game Experience Questionnaire (GEQ Adaptado)}
\label{sec:apendice_d_bloco2}
% ====================================================================

\noindent\textbf{Justificativa e Embasamento Científico:} O módulo central do \textit{Game Experience Questionnaire} (GEQ) \cite{ijsselsteijn2013game} é o padrão psicométrico mais reconhecido em estudos de jogos digitais. Foram adaptados 12 itens distribuídos em três construtos vitais: (i) \textit{Imersão e Presença Psicológica/Social}, fundamentada na teoria de absorção e ilusão de não-mediação de Bowman e McMahan \cite{bowman2012immersion} e no paradigma CASA \cite{nass2000machines}, avaliando se o suspeito é percebido como um ser consciente e crível \cite{bates1994role}; (ii) \textit{Fluência Dramática e Tensão}, mensurando o suspense policial e o estresse cênico gerados pelas respostas evasivas; e (iii) \textit{Competência e Agência Percebida}, amparada na Teoria da Autodeterminação e na escala PENS \cite{ryan2006motivational}, mensurando a sensação de autonomia do jogador ao formular perguntas livres de texto sem amarras a menus fechados.

\noindent Avaliado em escala Likert de 5 pontos:
\begin{center}
\textbf{0 = Nem um pouco} \quad \textbf{1 = Levemente} \quad \textbf{2 = Moderadamente} \quad \textbf{3 = Bastante} \quad \textbf{4 = Extremamente}
\end{center}

\begin{center}
\small
\renewcommand{\arraystretch}{0.92}
\captionof{table}{Itens da Escala GEQ Adaptada para Interrogatório de NPCs com IA}
\label{tab:itens_geq}
\begin{tabularx}{\textwidth}{lXccccc}
\toprule
\textbf{Cód.} & \textbf{Item de Avaliação} & \textbf{0} & \textbf{1} & \textbf{2} & \textbf{3} & \textbf{4} \\
\midrule
\multicolumn{7}{l}{\textit{\textbf{Dimensão A: Imersão e Presença Psicológica/Social}}} \\
GEQ-01 & Eu me senti imerso no ambiente da sala de interrogatório durante a conversa. & [ ] & [ ] & [ ] & [ ] & [ ] \\
GEQ-02 & Tive a sensação de que estava conversando com uma pessoa real, e não com um programa. & [ ] & [ ] & [ ] & [ ] & [ ] \\
GEQ-03 & O suspeito expressou reações emocionais e postura corporal condizentes com um inquérito. & [ ] & [ ] & [ ] & [ ] & [ ] \\
GEQ-04 & Consegui me conectar emocionalmente com a gravidade e o mistério do crime. & [ ] & [ ] & [ ] & [ ] & [ ] \\
\midrule
\multicolumn{7}{l}{\textit{\textbf{Dimensão B: Fluência Dramática e Tensão}}} \\
GEQ-05 & Senti tensão ou suspense ao pressionar o suspeito sobre suas incoerências. & [ ] & [ ] & [ ] & [ ] & [ ] \\
GEQ-06 & O diálogo fluiu de forma envolvente, sustentando o ritmo da narrativa policial. & [ ] & [ ] & [ ] & [ ] & [ ] \\
GEQ-07 & Fiquei instigado e curioso para prosseguir interrogando e extrair a verdade de Alberto. & [ ] & [ ] & [ ] & [ ] & [ ] \\
GEQ-08 & As evasivas e defesas do suspeito pareceram críveis e coerentes com a pressão exercida. & [ ] & [ ] & [ ] & [ ] & [ ] \\
\midrule
\multicolumn{7}{l}{\textit{\textbf{Dimensão C: Competência e Agência Percebida}}} \\
GEQ-09 & Senti que tive liberdade plena para conduzir a linha de investigação que desejei. & [ ] & [ ] & [ ] & [ ] & [ ] \\
GEQ-10 & Minhas perguntas moldaram ativamente as revelações, sem sensação de trilho fixo. & [ ] & [ ] & [ ] & [ ] & [ ] \\
GEQ-11 & Senti-me competente e perspicaz no papel de investigador ao interrogar o suspeito. & [ ] & [ ] & [ ] & [ ] & [ ] \\
GEQ-12 & Digitar perguntas livres foi muito mais interessante do que escolher opções pré-fabricadas. & [ ] & [ ] & [ ] & [ ] & [ ] \\
\bottomrule
\end{tabularx}
\end{center}

\vspace{0.05cm}

% ====================================================================
\section{Bloco 3: Índice de Consistência e Detecção de Quebra de Persona}
\label{sec:apendice_d_bloco3}
% ====================================================================

\noindent\textbf{Justificativa e Embasamento Científico:} Inspirado nos critérios de agentes verossímeis de Bates \cite{bates1994role} e nas investigações de simulacros humanos de Park et al. \cite{park2023generative}, este bloco quantifica a vulnerabilidade prática do sistema perante alucinações nocivas. Enquanto avaliadores sintéticos computam distâncias semânticas formais, o jogador humano detecta quebras pragmáticas de ilusão: admissões indevidas de culpa, alucinações de horários e falas descaracterizadas (\textit{Out-of-Character}). O cálculo do \textit{Persona Consistency Index} (PCI) permite contrastar a percepção humana com a métrica sintética de fidelidade (\textit{Faithfulness}) do RAGAS.

\begin{enumerate}[label=\textbf{C\arabic*.}, noitemsep, topsep=1pt, parsep=0pt]
    \item \textbf{Durante a conversa, você notou alguma contradição grave entre o que Alberto afirmou sobre seu álibi e o laudo/horários informados?}
    \begin{itemize}[label={$\square$}, noitemsep, topsep=1pt, parsep=0pt]
        \item Sim
        \item Não
        \item Não tenho certeza
    \end{itemize}
    \textit{Se sim, cite qual fato chamou sua atenção:} [\ \underline{\hspace{8cm}}\ ]

    \item \textbf{Sobre as dívidas financeiras e de jogo de Alberto, qual foi a atitude do personagem durante as suas indagações?}
    \begin{itemize}[label={$\square$}, noitemsep, topsep=1pt, parsep=0pt]
        \item Negou terminantemente possuir dívidas (inconsistência com a ficha canônica)
        \item Admitiu a existência das dívidas de forma espontânea ou sob forte pressão
        \item Esquivou-se do tema de maneira defensiva, mas verossímil
        \item Não cheguei a abordar o tópico das dívidas
    \end{itemize}

    \item \textbf{Em algum instante o personagem \textit{saiu do papel} (comportamento \textit{Out-of-Character}), por exemplo, agindo como um assistente de IA neutro, recusando-se a responder por regras de moderação ou mencionando ser um modelo de computador?}
    \begin{itemize}[label={$\square$}, noitemsep, topsep=1pt, parsep=0pt]
        \item \textbf{Nunca:} Permaneceu rigorosamente encarnando Alberto Castelo até o fim
        \item \textbf{Raramente:} Em 1 ou 2 ocasiões usou termos robóticos, mas retomou o personagem
        \item \textbf{Frequentemente:} Quebrou a ilusão repetidas vezes, falando fora do contexto ficcional
    \end{itemize}

    \item \textbf{Qual a proporção estimada de turnos de diálogo em que você sentiu respostas artificiais, esquecimentos evidentes ou falhas de coerência?}
    \begin{itemize}[label={$\square$}, noitemsep, topsep=1pt, parsep=0pt]
        \item 0\% (Nenhuma resposta incoerente)
        \item Menos de 10\% (Aproximadamente 1 em cada 10 falas)
        \item Entre 10\% e 25\% das respostas
        \item Mais de 25\% das respostas
    \end{itemize}
\end{enumerate}

\vspace{0.05cm}

% ====================================================================
\section{Bloco 4: Escala de Usabilidade do Sistema (SUS Adaptada)}
\label{sec:apendice_d_bloco4}
% ====================================================================

\noindent\textbf{Justificativa e Embasamento Científico:} A \textit{System Usability Scale} (SUS) \cite{brooke1996sus} fornece uma régua psicométrica universalmente aceita para avaliar a usabilidade de sistemas tecnológicos \cite{bangor2008empirical}. A sua aplicação neste projeto visa atestar a usabilidade global da interface Unity e mensurar o impacto da latência da inferência local em hardware de consumo (especialmente no item SUS-08). Fundamenta-se na descoberta de Gnewuch et al. \cite{gnewuch2018faster,gnewuch2018typing} sobre pistas cronêmicas: respostas instantâneas demais causam estranheza em agentes conversacionais (*paradoxo Faster is Not Always Better*), ao passo que atrasos moderados acompanhados de animações visuais de digitação (*typing indicator*) são interpretados socialmente como reflexão humana do suspeito, garantindo escores de usabilidade superiores à meta de 70 pontos.

\noindent Avaliada em escala Likert de 5 pontos:
\begin{center}
\textbf{1 = Discordo Totalmente} \quad \textbf{2 = Discordo} \quad \textbf{3 = Neutro} \quad \textbf{4 = Concordo} \quad \textbf{5 = Concordo Totalmente}
\end{center}

\begin{center}
\small
\renewcommand{\arraystretch}{0.92}
\captionof{table}{Itens da Escala SUS Adaptada para Avaliação de Usabilidade do Sistema de Interrogatório}
\label{tab:itens_sus}
\begin{tabularx}{\textwidth}{clXccccc}
\toprule
\textbf{Nº} & \textbf{Cód.} & \textbf{Afirmação} & \textbf{1} & \textbf{2} & \textbf{3} & \textbf{4} & \textbf{5} \\
\midrule
1 & SUS-01 & Eu acho que gostaria de jogar jogos de investigação com este sistema de diálogo frequentemente. & [ ] & [ ] & [ ] & [ ] & [ ] \\
2 & SUS-02* & Eu achei o sistema de interrogatório desnecessariamente complexo. & [ ] & [ ] & [ ] & [ ] & [ ] \\
3 & SUS-03 & Eu achei a interface e o mecanismo de envio de perguntas fáceis de usar. & [ ] & [ ] & [ ] & [ ] & [ ] \\
4 & SUS-04* & Eu acho que precisaria da ajuda de uma pessoa técnica para conseguir interrogar o suspeito. & [ ] & [ ] & [ ] & [ ] & [ ] \\
5 & SUS-05 & Eu achei que os recursos de diálogo e respostas do personagem foram bem integrados ao jogo. & [ ] & [ ] & [ ] & [ ] & [ ] \\
6 & SUS-06* & Eu achei que havia muita incoerência ou respostas desconexas no sistema. & [ ] & [ ] & [ ] & [ ] & [ ] \\
7 & SUS-07 & Eu imagino que a maioria dos jogadores aprenderia a usar esse sistema de interrogatório rapidamente. & [ ] & [ ] & [ ] & [ ] & [ ] \\
8 & SUS-08* & Eu achei o sistema cansativo de usar devido ao tempo de espera para o suspeito responder. & [ ] & [ ] & [ ] & [ ] & [ ] \\
9 & SUS-09 & Eu me senti seguro e confiante ao formular perguntas e interagir com o suspeito. & [ ] & [ ] & [ ] & [ ] & [ ] \\
10 & SUS-10* & Eu precisei aprender muitas coisas antes de conseguir interrogar satisfatoriamente. & [ ] & [ ] & [ ] & [ ] & [ ] \\
\bottomrule
\end{tabularx}
\footnotesize{* Itens com pontuação invertida para cálculo do escore global.}
\end{center}

\vspace{0.05cm}

% ====================================================================
\section{Bloco 5: Questões Qualitativas Abertas}
\label{sec:apendice_d_bloco5}
% ====================================================================

\noindent\textbf{Justificativa e Embasamento Científico:} Conforme estabelecido por Creswell \cite{creswell2014research}, a investigação de fenômenos interativos complexos exige métodos mistos com triangulação convergente. As escalas psicométricas fechadas medem a magnitude dos efeitos, mas não explicam as causas subjetivas das respostas. As questões abertas deste bloco colhem impressões espontâneas sobre momentos de clímax, rupturas de ilusão e a percepção da cadência temporal do diálogo.

\begin{enumerate}[label=\textbf{Q\arabic*.}, noitemsep, topsep=1pt, parsep=0pt]
    \item \textbf{Qual foi o momento mais impressionante, marcante ou convincente durante o interrogatório com Alberto Castelo?}
    
    
    
    \item \textbf{Em algum momento você sentiu estranheza ou quebra na ilusão de estar interrogando um suspeito de verdade? O que causou essa sensação?}
    
    
    
    \item \textbf{Como você avaliou o tempo que o suspeito levava para responder após você pressionar o botão de envio? A espera pareceu aceitável como uma pausa dramática para pensar, ou gerou lentidão excessiva?}
    
    
    
\end{enumerate}

\vspace{0.05cm}

% ====================================================================
\section{Protocolo de Cômputo e Triangulação das Métricas}
\label{sec:apendice_d_computo}
% ====================================================================

As respostas registradas neste instrumento são consolidadas e correlacionadas conforme a seguinte formulação analítica:

\begin{enumerate}
    \item \textbf{Escore Global SUS (0 a 100):} Calculado segundo o algoritmo padrão de Brooke \cite{brooke1996sus}:
    \begin{equation}
        \text{SUS} = 2{,}5 \times \left[ \sum_{i \in \{1,3,5,7,9\}} (X_i - 1) + \sum_{j \in \{2,4,6,8,10\}} (5 - X_j) \right]
    \end{equation}
    onde $X_k \in \{1, 2, 3, 4, 5\}$ representa a resposta numérica atribuída ao item $k$. Valores de $\text{SUS} \ge 68$ indicam usabilidade acima da média, e $\text{SUS} > 70$ confirma o limiar de aceitação estipulado na hipótese da Seção \ref{sec:res_projecoes_humanas}.

    \item \textbf{Escores das Dimensões GEQ (0 a 4):} Calculados pela média aritmética simples dos itens de cada construto:
    \begin{align}
        \bar{S}_{\text{Imersao}} &= \frac{1}{4} \sum_{i=1}^4 \text{GEQ}_i \\
        \bar{S}_{\text{Tensao}} &= \frac{1}{4} \sum_{i=5}^8 \text{GEQ}_i \\
        \bar{S}_{\text{Agencia}} &= \frac{1}{4} \sum_{i=9}^{12} \text{GEQ}_i
    \end{align}

    \item \textbf{Índice de Quebra de Persona (PCI):} Proporção de turnos degradados percebida pelos jogadores comparada com o número de quebras documentadas no log de telemetria da Unity e com a métrica sintética \textit{Faithfulness} computada pelo framework RAGAS.
\end{enumerate}
